\chapter{Gastric Emptying Study}

\section*{Overview}
A gastric emptying study measures how quickly food leaves the stomach using a small amount of radiolabeled material and serial imaging. The exact approach, setting, and preparation depend on the indication, anatomy, and the patient's health.

\section*{Why It Is Done}
It evaluates suspected gastroparesis, rapid gastric emptying, unexplained nausea and vomiting, early fullness, bloating, or upper-abdominal symptoms.

\section*{How It Is Performed}
After an overnight fast, the patient eats a standardized meal containing a tiny tracer. A gamma camera obtains images over several hours while the patient remains otherwise relaxed.

\section*{Preparation}
Follow fasting and medicine instructions; diabetes medicines and medicines that change stomach motility may require special planning. Do not alter them without advice.

\section*{Risks and Limitations}
Radiation exposure is low, and allergic reactions are not expected from the tracer. Results can be affected by the meal, blood glucose, medicines, and failure to complete scheduled images.

\section*{Results and Interpretation}
Results report the percentage of meal remaining at specified times. Delayed or rapid emptying is interpreted with symptoms, diabetes status, medicines, and other gastrointestinal findings.

\section*{Aftercare and Recovery}
Most people return to normal activity after the study. Ask the team when to restart medicines that were withheld and follow its instructions about eating and drinking.

\section*{Contraindications and Precautions}
The team reviews pregnancy, diabetes and blood-glucose control, medicines that alter stomach motility, inability to eat the standardized meal, and recent gastrointestinal procedures. Medicines should be changed only according to the ordering team's plan.

\section*{When to Seek Medical Care}
Follow the procedure team's specific instructions. Seek urgent medical assessment for severe or worsening pain, uncontrolled bleeding, fever or signs of infection, trouble breathing, chest pain, fainting, new weakness or confusion, inability to urinate, or any rapidly worsening symptom. The appropriate warning signs depend on the procedure and the patient's medical condition.

\section*{References}
\begin{enumerate}
  \item MedlinePlus. Gastric Emptying Study. U.S. National Library of Medicine; 2025.
  \item Current Medical Diagnosis and Treatment. McGraw-Hill; 2025.
\end{enumerate}
\clearpage
